\chapter[Introducción]{Introducción}
\label{cap:intro}

\section{Problema}
\label{sec:problema}

% PROVISORIO (01/10/2026): Milesight se cita como ejemplo verificado. Cuando el profesor guía confirme que
% los equipos adquiridos por la Facultad son Milesight, «Un fabricante de este tipo de equipos, Milesight,»
% pasa a «El fabricante de los equipos adquiridos por la Facultad, Milesight,» (pendiente 17).
La Facultad de Ingeniería de la Universidad Católica de Temuco está incorporando dispositivos IoT a su
campus. Un fabricante de este tipo de equipos, Milesight, entrega su propia plataforma, Milesight IoT
Cloud, que acepta solo productos de su marca y no ofrece una API para llevar los datos a otros sistemas
\parencite{milesightfaq}. Con dos marcas hay dos plataformas, y la información queda repartida, sujeta
a los límites de cada plan y difícil de reunir para investigación.

\begin{quote}
\textbf{Problema central:} la organización que ya posee dispositivos IoT no logra convertir sus datos
en información utilizable, porque carece de un sistema integrado que los registre, reciba y unifique
sus datos, vigile su estado y los presente en un solo lugar.
\end{quote}

El análisis identifica cinco causas: \textbf{C1 formatos}, cada fabricante emite sus datos en un
formato distinto; \textbf{C2 hardware}, la plataforma viene atada al equipo; \textbf{C3 nube}, costo
por dispositivo y datos fuera de la institución; \textbf{C4 técnicos}, no hay personal dedicado a
operar una plataforma; y \textbf{C5 visibilidad}, el silencio de un dispositivo es indistinguible de
la ausencia de novedades. La figura~\ref{fig:arbol} muestra el árbol completo.

\begin{figure}[H]
\centering
\resizebox{\linewidth}{!}{%
\begin{tikzpicture}[font=\scriptsize,
  cuadro/.style={caja, text width=2.0cm, minimum height=1.5cm, align=center}]

% ---- EFECTOS ----
\node[cuadro] (e1) {\textbf{E1}\\ Inversión\\ inmovilizada};
\node[cuadro, right=3mm of e1] (e2) {\textbf{E2}\\ Datos\\ perdidos};
\node[cuadro, right=3mm of e2] (e3) {\textbf{E3}\\ Decisiones\\ sin dato};
\node[cuadro, right=3mm of e3] (e4) {\textbf{E4}\\ Fallas vistas\\ tarde};
\node[cuadro, right=3mm of e4] (e5) {\textbf{E5}\\ Dependencia\\ del proveedor};
\begin{scope}[on background layer]
\node[grupo, fit=(e1)(e5)] (gefectos) {};
\end{scope}
\node[etiqueta, left=2mm of gefectos] {EFECTOS};

% ---- PROBLEMA CENTRAL ----
\node[cajaacento, below=11mm of gefectos, text width=11.6cm, minimum height=1.15cm,
      fill=verdeacento!6, font=\scriptsize, align=center] (prob)
  {\textbf{Quien ya posee sensores no puede convertir sus datos en información utilizable:}\\
   falta una plataforma que unifique la recepción, la normalización y la gestión de dispositivos.};
\node[etiqueta, left=2mm of prob] {PROBLEMA\\CENTRAL};

% ---- CAUSAS (con apodo memorizable) ----
\node[cuadro, below=11mm of prob, xshift=-4.65cm] (c1) {\textbf{C1}\\ Formatos\\ dispares};
\node[cuadro, right=3mm of c1] (c2) {\textbf{C2}\\ Hardware\\ atado};
\node[cuadro, right=3mm of c2] (c3) {\textbf{C3}\\ Costo y\\ nube ajena};
\node[cuadro, right=3mm of c3] (c4) {\textbf{C4}\\ Sin personal\\ técnico};
\node[cuadro, right=3mm of c4] (c5) {\textbf{C5}\\ Sin\\ visibilidad};
\begin{scope}[on background layer]
\node[grupo, fit=(c1)(c5)] (gcausas) {};
\end{scope}
\node[etiqueta, left=2mm of gcausas] {CAUSAS};

% ---- flechas perfectamente verticales, ancladas al centro del problema ----
\draw[flechagruesa] (prob.south |- gcausas.north) -- node[nota, right=1mm] {provocan} (prob.south);
\draw[flechagruesa] (prob.north |- gefectos.south) -- node[nota, right=1mm] {produce} (prob.north);
\end{tikzpicture}}
\caption{Árbol de problemas. Las cinco causas —C1 formatos dispares, C2 hardware atado, C3 costo y
nube ajena, C4 sin personal técnico, C5 sin visibilidad— provocan el problema central, que produce
los cinco efectos.}
\label{fig:arbol}
\end{figure}


\textbf{Evidencia.} La literatura identifica la fragmentación de protocolos y formatos como uno de los
desafíos centrales del Internet de las Cosas \parencite{alfuqaha2015} y documenta que la falta de
interoperabilidad y la dependencia de proveedor impiden su adopción a gran escala \parencite{noura2019},
lo que sustenta C1, C2 y C3. La evidencia directa de la Facultad ---necesidades, incidentes y uso
previsto de los datos--- se está levantando por entrevista (sección~\ref{sec:retro}).

\textbf{Organización destinataria y alcance.} La Universidad Católica de Temuco es la dueña del
problema; el caso de validación es el Campus Dr.\ Luis Rivas del Canto, sede de la Facultad de
Ingeniería. El trabajo abarca el sistema integrado, desde que un dispositivo IoT se registra hasta que
sus datos se muestran en el dashboard, incluido el diseño de la red por la que esos datos viajan desde
el sitio hasta el servidor. Los dispositivos y su radio no son parte de este trabajo. Quedan fuera la
fabricación de dispositivos y la analítica predictiva.

\section{Estado del arte}
\label{sec:estado-arte}

Hoy existen sistemas integrados de código abierto para registrar, gestionar y monitorear dispositivos
IoT, como OpenRemote, Magistrala y ThingsBoard; este último cambió de licencia en septiembre de 2026.
Existen además soluciones atadas a su propio hardware o a su servicio, como Milesight IoT Cloud o la
chilena Energetika, y plataformas comerciales en la nube, como Ubidots y TagoIO.

OpenRemote, bajo licencia AGPL~3.0 \parencite{openremotelicencia}, es el más completo de los abiertos:
paneles, reglas, alarmas, notificaciones por correo y al celular, llamadas a sistemas externos y mapa,
e integra redes LoRaWAN a través de ChirpStack; su instalación central exige además un servidor de
identidad, Keycloak \parencite{openremotedocs}. Magistrala libera su núcleo bajo Apache~2.0
\parencite{magistralalicencia}, pero deja los paneles, las alarmas, el motor de reglas y los reportes
fuera de su edición comunitaria \parencite{magistraladocs}. ThingsBoard, el más difundido, distribuye
desde su versión~4.4 su código bajo la \emph{Business Source License}, que permite el uso gratuito
con condiciones ---para una institución educativa, hasta 1\,000 dispositivos--- \parencite{thingsboardbusl};
la versión~4.3 sigue en Apache~2.0 \parencite{thingsboardlicencia43}. Su guía de instalación pide
4\,GB de memoria para pruebas y 8\,GB para una producción pequeña \parencite{thingsboardinstalacion}.

Milesight IoT Cloud acepta solo productos Milesight conectados a través de sus propias pasarelas, no
ofrece una API hacia otras plataformas y permite descargar los datos en CSV de a tres meses
\parencite{milesightfaq}; su plan gratuito admite 10 dispositivos y 2 paneles \parencite{milesightcloud}.
Energetika ofrece en Chile una plataforma propia, desarrollada y operada por la empresa, que trabaja
con los sensores que el cliente ya tiene \parencite{energetikaiotify}.

Ubidots y TagoIO aceptan equipos que no están en su catálogo \parencite{ubidotsdispositivos,tagodecodificador},
pero alojan los datos en servidores del proveedor, en Estados Unidos o Europa
\parencite{ubidotsseguridad,tagoapi}, con planes gratuitos de 3 y 5 dispositivos; el de Ubidots, solo
para uso no comercial y con un mes de retención \parencite{ubidotsstem,tagoprecios}.

Cada grupo resuelve una parte del problema y deja otra: las soluciones del fabricante atan los datos
a su hardware; las nubes comerciales limitan los dispositivos y guardan los datos fuera de la
institución; y las abiertas, que se instalan en un servidor propio, exigen componentes adicionales,
reservan funciones a su edición de pago o condicionan su uso. Este trabajo se ubica en el grupo de
código abierto, como prototipo: no compite en amplitud de funciones con OpenRemote o ThingsBoard.
Busca otra cosa: que el personal no técnico de una facultad registre y vigile sus dispositivos y deje
sus datos listos para investigación, exportables en formato abierto, con independencia del fabricante
y en la propia institución.

\section{Objetivos}
\label{sec:objetivos}
\subsection{Objetivo general}
\begin{itemize}
\item Desarrollar una plataforma web para el registro, monitoreo y gestión de datos de sensores,
que reciba las tramas de distintos sensores a través de TCP/IP, las normalice a un modelo interno
único, almacene las mediciones como series de tiempo, supervise el estado de cada dispositivo y
presente la información en un panel de control, validada sobre un caso de uso real en el Campus
Dr. Luis Rivas del Canto de la Universidad Católica de Temuco al término del semestre académico 2026.
\end{itemize}

\subsection{Objetivos específicos}
\begin{itemize}
\item \textbf{OE1} — Caracterizar la situación de gestión de datos de sensores en la organización
destinataria, mediante revisión de literatura, análisis de la oferta existente y al menos dos entrevistas
semiestructuradas documentadas (participantes, necesidades, prioridades), para validar un catálogo
de requerimientos trazable a las causas del problema, al 30 de octubre de 2026.
\item \textbf{OE2} — Evaluar y seleccionar el mecanismo de transporte y el formato de trama, y
especificar el contrato de entrada con sus adaptadores y política de cuarentena, mediante matrices
de criterios ponderados, para obtener el documento de especificación de la capa de recepción, al 2
de octubre de 2026.
\item \textbf{OE3} — Diseñar la arquitectura de la plataforma, su modelo de datos y su despliegue
con costo cuantificado, mediante la fase \emph{Design} de PPDIOO y la comparación de al menos tres
alternativas de alojamiento, para obtener los diagramas, el modelo y la recomendación justificada,
al 2 de octubre de 2026.
\item \textbf{OE4} — Implementar la plataforma, mediante desarrollo incremental, para obtener
primero el \textbf{núcleo mínimo demostrable} —recepción por MQTT, normalización, supervisión de
caída y consulta básica— con un sensor real de extremo a extremo al 30 de octubre de 2026, y luego
la entrada HTTP, las alertas y el panel de control al 4 de diciembre de 2026.
\item \textbf{OE5} — Validar la plataforma en un servidor de recursos modestos, mediante pruebas
funcionales, de carga y de recuperación, para obtener un informe de validación contra los umbrales
de los requerimientos no funcionales —al menos 15 casos de prueba con el 90\,\% aprobados; 50
mensajes por segundo sin pérdida; latencia de recepción a consulta de 2 segundos o menos en el
percentil 95; caída detectada en a lo más dos intervalos de reporte; recuperación automática sin
perder mensajes confirmados; disponibilidad de la recepción del 99\,\% o más en una ventana continua
de 72 horas; memoria total de 2\,GB o menos—, al 4 de diciembre de 2026.
\end{itemize}

El cuadro~\ref{tab:trazabilidad} resume el producto verificable de cada objetivo y su trazabilidad;
los hitos y metas numéricas están en el repositorio y en el capítulo~\ref{cap:metodologia}.

\begin{table}[H]
\centering\footnotesize
\begin{tablaseg}
\begin{tabular}{@{}L{2.9cm}L{3.5cm}L{1.9cm}L{1.3cm}L{2.2cm}C{0.9cm}@{}}
\toprule
\textbf{Objetivo} & \textbf{Producto verificable} & \textbf{Causas} & \textbf{Efectos} & \textbf{Fase PPDIOO} & \textbf{Etapa} \\
\midrule
OE1 Caracterizar & Catálogo de requerimientos con fuente por ítem (anexo~\ref{anx:requerimientos}); registros de entrevista & C4 (diagnostica C1--C5) & E1, E3 & Prepare & 1, 3 \\
OE2 Contrato de recepción & Especificación del contrato: tópico, SenML, adaptadores, cuarentena (cap.~\ref{cap:especificacion}) & \textbf{C1} & E2, E5 & Plan, Design & 2 \\
OE3 Arquitectura y despliegue & Vistas C4, ER y UML; recomendación de alojamiento con costo (cap.~\ref{cap:diseno}) & C2, C3, C5 & E1, E4, E5 & Design, Plan & 2 \\
OE4 Implementar & Prototipo en repositorio: núcleo mínimo y luego segundo nivel & C1, C2, C4 & E1--E3 & Implement & 3, 4 \\
OE5 Validar & Informe de validación contra los umbrales & C5 & E4 & Operate, Optimize & 4 \\
\bottomrule
\end{tabular}
\end{tablaseg}
\caption{Matriz de trazabilidad: producto verificable de cada objetivo, causas y efectos que ataca,
fase y etapa. Las cinco causas y los cinco efectos quedan cubiertos por al menos un objetivo.}
\label{tab:trazabilidad}
\end{table}

\section{Autoría y licenciamiento}

Este trabajo, su documentación y el código son de \textbf{autoría de Vicente Ricardo Ramírez Oñate};
la universidad es la organización destinataria del producto, no la titular de la obra. La licencia
de publicación se decidirá antes de liberar el resultado, evaluando opciones que permitan el uso
abierto e impidan el aprovechamiento comercial por terceros sin acuerdo con el autor.

\section{Estructura del documento}

Antecedentes técnicos (cap.~\ref{cap:antecedentes}); requerimientos y evaluaciones
(cap.~\ref{cap:especificacion}); diseño desde varias vistas, nuevo en esta etapa
(cap.~\ref{cap:diseno}); metodología, seguimiento y retroalimentación incorporada
(cap.~\ref{cap:metodologia}); anexos con el catálogo, el backlog, las evaluaciones, figuras
complementarias y el camino descartado.
